\documentclass{article}
\usepackage[T1]{fontenc}
\usepackage{lmodern}
\usepackage{graphicx}
\usepackage{amsmath,amssymb}
\usepackage[a4paper,margin=2cm]{geometry}
\usepackage[absolute,overlay]{textpos}
\setlength{\TPHorizModule}{1mm}
\setlength{\TPVertModule}{1mm}
\usepackage[indonesian]{babel}
\usepackage{hyperref}
\setlength{\emergencystretch}{2em}
% unit: Halaman 11

\begin{document}

% === Halaman 11 (python/native) ===

Karena bobot total tersisa = 0, maka solusi persoalan superincreasing 

knapsack ditemukan. Barisan bobot yang dimasukkan ke dalam 

knapsack adalah  


\{2, 3, – , 13, – , 52\}  


 sehingga

    70 = (1  2) + (1  3) + (0  6) + (1  13) + 


     (0  27) + (1  52)

 Dengan kata lain, plainteks dari kriptogram 70 adalah


 110101.

11

\end{document}
